\begin{tikzpicture}[x=1cm,y=1cm]
\node[note,anchor=east] at (0,2.1) {购买库 WAL};
\draw[arr] (.1,2.1)--(12,2.1);
\foreach \x/\t in {1/B0,4/A,8/R} {
  \fill[Blue] (\x,2.1) circle(2.5pt);
  \node[note,above] at (\x,2.2) {\t};
}
\node[note] at (10,2.65) {原时间线：已有 P1/C1/O1/EV1};
\draw[arr,draw=Gold] (1,2.1)--(3,.6)--(12,.6);
\node[warn,text width=4.2cm] at (8.4,.6)
  {从 B0 开始的新时间线\\本地尚无原 S/C1};
\node[note,above] at (3.6,1.2) {选择旧停止点};
\draw[lane] (.1,-.6)--(12,-.6);
\node[note,anchor=east] at (0,-1.3) {平台账本};
\draw[arr] (.1,-1.3)--(12,-1.3);
\fill[green!35!black] (6,-1.3) circle(2.5pt);
\node[note,above] at (6,-1.2) {X1，余额 7000};
\node[note,anchor=east] at (0,-2.5) {履约保存域};
\draw[arr] (.1,-2.5)--(12,-2.5);
\fill[green!35!black] (10,-2.5) circle(2.5pt);
\node[note,above] at (10,-2.4) {原 I/B，F=501};
\draw[lane] (6,1.8)--(6,-1);
\draw[lane] (10,1.8)--(10,-2.2);
\node[note,text width=12cm] at (6,-3.3)
  {只有购买库选择旧状态；时间线区分 WAL 分支，跨域关联由业务证据恢复。};
\end{tikzpicture}
